\section{\journal{Results}}
\begin{journaladdition}
\noindent\textit{Journal draft: experiments are pending. All tables and figures below are planned outputs, not findings. The previous evaluation and leaderboard results have been removed because they used a different protocol.}

\subsection{Independent-Test Comparison}
\begin{itemize}
  \item Report the final cohort, exclusions, and evaluable structures/lesions.
  \item Compare baselines and the full model with and without IO on the frozen test partition (Table~\ref{tab:journal_baselines}).
  \item Add paired patient-level effects and uncertainty; describe alignment, preservation, folding, and runtime without selecting a winner from test results.
\end{itemize}
\end{journaladdition}
\begin{table}[htbp]
\journalcolor
\centering
\caption{\journal{Planned independent-test comparison. All entries are pending; ``pending'' is not a numerical result. Report patient-level uncertainty and paired effects in the final table/text. Per-lesion and regional endpoints are detailed separately.}}
\label{tab:journal_baselines}
\scriptsize
\setlength{\tabcolsep}{3pt}
\begin{tabular}{lccccc}
\toprule
Method & Dice & HD95 & MTV/TLG & NDV & Runtime \\
\midrule
Initial & pending & pending & pending & pending & pending \\
Affine & pending & pending & pending & pending & pending \\
NiftyReg & pending & pending & pending & pending & pending \\
ConvexAdam & pending & pending & pending & pending & pending \\
Full model, no IO & pending & pending & pending & pending & pending \\
Full model + IO & pending & pending & pending & pending & pending \\
\bottomrule
\end{tabular}
\end{table}


\begin{journaladdition}
\subsection{Separate and Combined Effects of PET Preservation and Rigidity}
\label{sec:ablation}
\begin{itemize}
  \item Fill the matched-training $2\times2$ comparison and the four PET-term removal conditions (Table~\ref{tab:journal_ablations}).
  \item Report aggregate versus per-lesion changes; check whether aggregate cancellation hides large individual errors.
  \item Show signed, absolute, and upper-tail errors by lesion size and bone association (Fig.~\ref{fig:journal_lesions}).
\end{itemize}
\end{journaladdition}
\begin{table}[htbp]
\journalcolor
\centering
\caption{\journal{Planned matched-training ablations. Accuracy and general regularization terms are retained. The endpoint groups are alignment, aggregate/per-lesion PET preservation, and deformation quality. Results and patient-level paired effects are pending; any additional IO-only ablations must be reported separately.}}
\label{tab:journal_ablations}
\small
\setlength{\tabcolsep}{4pt}
\begin{tabular}{lccc}
\toprule
Condition & PET terms & Rigidity & Results \\
\midrule
Neither & None & Off & Pending \\
PET only & All four & Off & Pending \\
Rigidity only & None & On & Pending \\
Both (full) & All four & On & Pending \\
\midrule
Full minus $\mathcal{L}_{\mathrm{MTV}}$ & Other three & On & Pending \\
Full minus $\mathcal{L}_{\mathrm{MTV\text{-}J}}$ & Other three & On & Pending \\
Full minus $\mathcal{L}_{\mathrm{jac}}$ & Other three & On & Pending \\
Full minus $\mathcal{L}_{\mathrm{TLG}}$ & Other three & On & Pending \\
\bottomrule
\end{tabular}
\end{table}

\begin{figure}[htbp]
\journalcolor
\centering
\figplaceholder{0.18}{Planned panels: aggregate versus per-lesion MTV/TLG error; signed and absolute error distributions; upper-tail summaries by lesion size and bone association.}
\caption{\journal{Planned lesion-preservation analysis for the matched PET and rigidity ablations. Compare each moving scan before and after warping, with patient-level uncertainty. State subgroup counts and the predeclared upper-tail statistic. No experimental data are shown.}}
\label{fig:journal_lesions}
\end{figure}


\begin{journaladdition}
\subsection{Rigidity Comparison and Boundary Effects}
\begin{itemize}
  \item Compare each rigidity method over its own weight range: bone/organ alignment, distance distortion, regional folding, and secondary BRE.
  \item Relate stencil leakage and erosion retention to the observed outcomes without assuming a mechanism in advance (Fig.~\ref{fig:journal_rigidity}).
  \item Include representative bone-boundary examples under a documented case-selection rule.
\end{itemize}
\end{journaladdition}
\begin{figure}[htbp]
\journalcolor
\centering
\figplaceholder{0.18}{Planned panels: bone/organ alignment versus physical distance distortion over method-specific weights; folding inside bone and in the outer band; stencil leakage and erosion retention; representative boundary maps.}
\caption{\journal{Planned comparison of no rigidity, Staring, Staring with erosion, and Kabsch. Report regional folding alongside alignment and distance distortion, and show boundary-mechanism measurements separately. BRE is secondary. No experimental data are shown.}}
\label{fig:journal_rigidity}
\end{figure}


\begin{journaladdition}
\subsection{Accuracy--Preservation Trade-off}
\begin{itemize}
  \item Show the predeclared PET-weight settings and their patient-level uncertainty (Fig.~\ref{fig:journal_tradeoff}).
  \item Identify the validation-selected operating point and report its frozen test performance.
  \item Discuss whether preservation changes alignment or folding, including heterogeneous lesion responses.
\end{itemize}
\end{journaladdition}
\begin{figure}[htbp]
\journalcolor
\centering
\figplaceholder{0.18}{Planned panels: alignment versus per-lesion MTV and TLG preservation across a small PET-weight sweep, including zero and the current setting; corresponding deformation quality.}
\caption{\journal{Planned accuracy--preservation trade-off. Mark the operating point selected on validation patients and distinguish validation curves from frozen-setting test results. Include patient-level uncertainty and identify whether settings change matched training or IO only. No experimental data are shown.}}
\label{fig:journal_tradeoff}
\end{figure}


\begin{journaladdition}
\subsection{Optional Supporting Results}
\todo{Include gradient analysis and/or external-cohort results only if completed; otherwise remove this subsection. Do not imply external validation from surrogate measurements alone.}
\end{journaladdition}
